\documentclass[11pt,a4paper]{article}

% ─────────────────────────────────────────────────────────────
% Packages
% ─────────────────────────────────────────────────────────────
\usepackage[utf8]{inputenc}
\usepackage[T1]{fontenc}
\usepackage[french]{babel}
\usepackage{lmodern}
\usepackage{geometry}
\usepackage{amsmath,amssymb,amsfonts}
\usepackage{bm}
\usepackage{graphicx}
\usepackage{booktabs}
\usepackage{hyperref}
\usepackage{physics}
\usepackage{caption}
\usepackage{subcaption}
\usepackage{xcolor}

\geometry{margin=2.5cm}

\hypersetup{
    colorlinks=true,
    linkcolor=blue,
    citecolor=blue,
    urlcolor=blue
}

% ─────────────────────────────────────────────────────────────
% Macros
% ─────────────────────────────────────────────────────────────
\newcommand{\bup}{\mathrm{BuP}}
\newcommand{\ent}{\mathrm{ent}}
\newcommand{\matter}{\mathrm{matter}}
\newcommand{\eff}{\mathrm{eff}}
\newcommand{\loc}{\mathrm{loc}}
\newcommand{\proxy}{\mathrm{proxy}}
\newcommand{\Omat}{\mathcal{O}_{\rm matter}}
\newcommand{\lc}{\lambda_c}

% ─────────────────────────────────────────────────────────────
% Title
% ─────────────────────────────────────────────────────────────
\title{
\textbf{Équation pure BuP : matière, localité et seuil critique d'intrication}\\[0.5em]
\large Matière et géométrie comme phases émergentes du réseau d'information mutuelle
}

\author{
Farid Hamdad\\
\small Independent Researcher\\
\small Bottom-Up Quantum Gravity Program
}

\date{2026}

\begin{document}

\maketitle

% ─────────────────────────────────────────────────────────────
\begin{abstract}
% ─────────────────────────────────────────────────────────────

Les travaux précédents du programme Bottom-Up Quantum Gravity (BuP) ont proposé que la distance, la métrique effective, la dimension spectrale et le couplage gravitationnel émergent du réseau d'information mutuelle
\[
W_{ij}=I(i:j)
\]
associé à un état quantique global. Paper~5 a formulé une équation d'Einstein effective dans laquelle la géométrie émergente répond aux variations locales de la dimension spectrale. Cependant, cette formulation conservait encore un tenseur énergie-impulsion de matière \(T_{\mu\nu}^{\rm matter}\) introduit comme source effective.

Le présent papier explore l'étape suivante : une formulation pure BuP dans laquelle la matière n'est plus postulée, mais apparaît comme une excitation localisée du réseau d'intrication. Nous proposons une équation fondamentale conjecturale
\[
\frac{\delta S_{\rm BuP}[W]}{\delta W_{ij}}=0,
\]
dont les limites continues doivent produire à la fois la géométrie émergente et les sources matérielles effectives.

Nous testons numériquement cette idée en injectant une perturbation localisée
\[
W_{ij}=W_{ij}^{(0)}+\delta W_{ij}^{\rm loc}
\]
dans des graphes d'information mutuelle, puis en comparant le proxy de matière
\[
T_{ij}^{\rm matter,proxy}=(\delta W_{ij})^2
\]
à la réponse de courbure induite
\[
|\delta G_{ij}^{\rm proxy}|.
\]
Nous définissons un candidat paramètre d'ordre
\[
\mathcal{O}_{\rm matter}(\lambda)
=
\rho_{\rm Spearman}
\left(
T_{ij}^{\rm matter,proxy},
|\delta G_{ij}^{\rm proxy}|
\right).
\]
Pour \(N=20\), cette observable change brutalement de signe autour de
\[
\lambda_c(N=20)\simeq0.599.
\]
Une analyse de taille finie préliminaire pour \(N=9,12,16,20\) indique que la transition n'est pas résolue aux petites tailles, mais émerge progressivement avec \(N\).

Nous interprétons ce résultat comme un seuil critique candidat de localité/matière : en dessous de \(\lambda_c\), les excitations localisées se comportent comme des sources géométriques effectives ; au-dessus, elles perdent leur cohérence géométrique locale. Cette étude ne constitue pas encore une dérivation complète de la matière, mais fournit un premier test numérique interne de l'idée que la matière peut émerger comme excitation stable d'une phase locale du réseau d'intrication.

\end{abstract}

\tableofcontents

% ─────────────────────────────────────────────────────────────
\section{Introduction}
% ─────────────────────────────────────────────────────────────

Le programme Bottom-Up Quantum Gravity (BuP) repose sur une hypothèse centrale : l'espace, la géométrie et la gravitation ne sont pas des structures fondamentales, mais émergent de la structure d'intrication d'un état quantique global. Le point de départ n'est donc pas une variété différentiable munie d'une métrique, mais un réseau d'information mutuelle
\[
W_{ij}=I(i:j),
\]
où \(I(i:j)\) mesure la corrélation informationnelle entre deux sous-systèmes quantiques.

Les travaux précédents ont progressivement développé cette idée. Une distance d'intrication a été définie à partir des poids \(W_{ij}\), puis utilisée pour reconstruire une métrique émergente \(g_{\mu\nu}^{\rm ent}\). Le Laplacien du graphe d'intrication a ensuite permis de définir une dimension spectrale locale \(d_s(x)\), et les variations de cette dimension ont été reliées à un couplage gravitationnel effectif
\[
G_{\rm eff}(d_s)
=
\frac{2}{d_s-1}G.
\]

Paper~5 a franchi une étape supplémentaire en proposant une équation d'Einstein effective émergente,
\[
G_{\mu\nu}[g^{\rm ent}]
=
8\pi G_{\rm eff}(d_s)T_{\mu\nu}^{\rm matter}
+
8\pi G\,T_{\mu\nu}^{\rm ent}[d_s].
\]
Dans cette équation, la métrique, la dimension spectrale et le couplage gravitationnel émergent du réseau d'intrication. Le terme
\[
T_{\mu\nu}^{\rm ent}[d_s]
\]
encode la réponse géométrique associée aux gradients et aux variations secondes de \(d_s\).

Cette formulation reste cependant partiellement effective. Le tenseur
\[
T_{\mu\nu}^{\rm matter}
\]
y est encore introduit comme une source de matière. Une formulation pleinement Bottom-Up devrait aller plus loin : la matière elle-même doit émerger du réseau d'intrication, au même titre que la géométrie.

Le présent papier propose cette étape. L'idée centrale est simple :

\[
\boxed{
\text{la matière n'est pas ajoutée au réseau ; elle est une excitation stable du réseau.}
}
\]

Autrement dit, une configuration de matière doit pouvoir être représentée comme une perturbation localisée
\[
W_{ij}
=
W_{ij}^{(0)}
+
\delta W_{ij}^{\rm loc}
\]
autour d'une configuration de vide \(W_{ij}^{(0)}\). Si cette perturbation est localisée, stable, transportable et géométriquement active, alors elle peut être interprétée comme une source matérielle effective après coarse-graining.

Le but de ce travail n'est pas encore de dériver explicitement les équations de Dirac, Maxwell ou Yang--Mills depuis le réseau \(W_{ij}\). L'objectif est plus modeste mais plus fondamental : formuler un cadre dans lequel la matière devient une propriété émergente du même substrat informationnel que la géométrie, puis tester numériquement si une excitation localisée de \(W_{ij}\) produit effectivement une réponse de courbure.

Nous proposons une équation pure BuP sous la forme variationnelle
\[
\boxed{
\frac{\delta S_{\rm BuP}[W]}{\delta W_{ij}}=0.
}
\]
Dans cette formulation, ni la métrique ni le tenseur de matière ne sont fondamentaux. Les deux doivent apparaître comme des limites continues ou coarse-grainées d'une unique dynamique informationnelle.

La contribution numérique principale de ce papier est la définition et la mesure d'un candidat paramètre d'ordre
\[
\mathcal{O}_{\rm matter}(\lambda)
=
\rho_{\rm Spearman}
\left(
T_{ij}^{\rm matter,proxy},
|\delta G_{ij}^{\rm proxy}|
\right),
\]
où
\[
T_{ij}^{\rm matter,proxy}=(\delta W_{ij})^2
\]
et
\[
G_{ij}^{\rm proxy}
=
-\kappa_{ij}
+
\frac12R_{\rm edge}.
\]
Cette observable mesure si une excitation localisée du réseau agit comme une source géométrique effective.

Nous montrons que, pour \(N=20\), \(\mathcal{O}_{\rm matter}(\lambda)\) change brutalement de signe autour de
\[
\lambda_c(N=20)\simeq0.599.
\]
Pour \(\lambda<\lambda_c\), l'excitation localisée produit une réponse de courbure positivement corrélée. Pour \(\lambda>\lambda_c\), la réponse devient anti-corrélée. Cette transition est indépendante du signe de la perturbation injectée, ce qui suggère qu'elle est une propriété globale du régime d'intrication, et non un artefact de l'excitation.

La structure du papier est la suivante. La Section~\ref{sec:paper5} rappelle le lien avec Paper~5. La Section~\ref{sec:pure-postulate} présente le postulat pur BuP. La Section~\ref{sec:action} introduit une action informationnelle conjecturale \(S_{\rm BuP}[W]\). La Section~\ref{sec:matter-excitations} décrit la matière comme excitation stable du réseau. La Section~\ref{sec:continuum-limits} discute les limites continues vers la géométrie et la matière. La Section~\ref{sec:numerical-program} présente le programme numérique. Les Sections~\ref{sec:results} et~\ref{sec:fss} exposent les résultats numériques et l'analyse de taille finie. La Section~\ref{sec:discussion} discute l'interprétation physique, notamment en termes de localité et de décohérence géométrique. La Section~\ref{sec:conclusion} conclut.

% ─────────────────────────────────────────────────────────────
\section{Lien avec Paper 5 : de la géométrie dynamique à la matière émergente}
\label{sec:paper5}
% ─────────────────────────────────────────────────────────────

Paper~5 a proposé une équation d'Einstein effective émergente :
\[
G_{\mu\nu}[g^{\rm ent}]
=
8\pi G_{\rm eff}(d_s)T_{\mu\nu}^{\rm matter}
+
8\pi G\,T_{\mu\nu}^{\rm ent}[d_s].
\]
Le tenseur d'intrication y est donné sous la forme
\[
T_{\mu\nu}^{\rm ent}[d_s]
=
\alpha
\left(
\nabla_\mu d_s\nabla_\nu d_s
-
\frac12 g_{\mu\nu}^{\rm ent}(\nabla d_s)^2
\right)
+
\beta
\left(
\nabla_\mu\nabla_\nu d_s
-
g_{\mu\nu}^{\rm ent}\Box d_s
\right).
\]

Un test numérique interne a montré qu'un proxy discret de courbure émergente
\[
G_{ij}^{\rm proxy}
=
-\kappa_{ij}
+
\frac12R_{\rm edge}
\]
est corrélé à des proxys spectraux construits à partir des variations locales de \(d_s\), notamment
\[
(\nabla d_s)^2
\qquad\text{et}\qquad
\Delta d_s.
\]
Ce résultat fournit une première cohérence numérique entre courbure émergente et tenseur spectral.

Le présent papier utilise ce résultat comme base. Si le graphe d'intrication permet déjà de reconstruire une réponse géométrique de type Einstein, alors on peut poser la question suivante :

\[
\boxed{
\text{une excitation localisée de } W_{ij} \text{ peut-elle jouer le rôle de source matérielle ?}
}
\]

Autrement dit, Paper~5 traite la partie
\[
W_{ij}
\longrightarrow
g_{\mu\nu}^{\rm ent},
d_s(x),
G_{\mu\nu}^{\rm ent}.
\]
Paper~6 explore la partie complémentaire :
\[
\delta W_{ij}^{\rm loc}
\longrightarrow
T_{\mu\nu}^{\rm matter}.
\]

La transition conceptuelle est donc :
\[
\text{géométrie émergente}
\quad\longrightarrow\quad
\text{matière émergente}.
\]

% ─────────────────────────────────────────────────────────────
\section{Postulat pur BuP}
\label{sec:pure-postulate}
% ─────────────────────────────────────────────────────────────

Le postulat pur BuP peut être formulé ainsi :
\[
\boxed{
\text{aucune géométrie, aucune matière et aucun champ ne sont fondamentaux.}
}
\]

Le seul objet fondamental est l'état quantique global et la structure d'intrication qu'il induit. Cette structure est représentée par la matrice
\[
W_{ij}=I(i:j).
\]

En relativité générale, l'équation
\[
G_{\mu\nu}=8\pi G T_{\mu\nu}
\]
sépare deux secteurs : la géométrie à gauche, la matière à droite. Dans une formulation pure BuP, cette séparation ne peut pas être fondamentale. Les deux membres doivent apparaître comme des descriptions effectives du même réseau d'intrication.

La chaîne d'émergence attendue est :
\[
W_{ij}
\quad\Longrightarrow\quad
\begin{cases}
g_{\mu\nu}^{\rm ent},\\
d_s(x),\\
G_{\rm eff}(d_s),\\
T_{\mu\nu}^{\rm matter},\\
\text{champs effectifs}.
\end{cases}
\]

Les trois premiers objets ont déjà été étudiés dans les travaux précédents. Le présent papier se concentre sur le quatrième : l'émergence du tenseur de matière.

On note
\[
W_{ij}^{(0)}
\]
une configuration de vide du réseau. Une excitation localisée est une perturbation
\[
W_{ij}
=
W_{ij}^{(0)}
+
\delta W_{ij}^{\rm loc}.
\]
Une telle excitation peut être interprétée comme matière si elle est localisée, stable, transportable et si elle produit une réponse géométrique locale.

Dans cette perspective, le tenseur énergie-impulsion n'est pas ajouté au réseau. Il devient une fonctionnelle coarse-grainée des excitations :
\[
T_{\mu\nu}^{\rm matter}
=
\mathcal{T}_{\mu\nu}[\delta W^{\rm loc}].
\]

L'équation pure BuP s'écrit donc sous la forme
\[
\frac{\delta S_{\rm BuP}[W]}{\delta W_{ij}}=0.
\]
Les équations de la géométrie et celles de la matière doivent émerger comme limites continues de cette équation unique.

% ─────────────────────────────────────────────────────────────
\section{Action informationnelle du réseau}
\label{sec:action}
% ─────────────────────────────────────────────────────────────

Pour donner un contenu dynamique à l'équation pure BuP, il faut spécifier une action effective définie sur le réseau d'intrication. Une forme minimale conjecturale est
\[
S_{\rm BuP}[W]
=
S_{\rm spec}[L]
+
S_{\rm curv}[W]
+
S_{\rm loc}[W]
+
S_{\rm topo}[W],
\]
où
\[
L=D-W,
\qquad
D_{ii}=\sum_j W_{ij}
\]
est le Laplacien du graphe.

Le terme spectral
\[
S_{\rm spec}[L]
=
{\rm Tr}\,f(L)
\]
contrôle le spectre du Laplacien, donc la diffusion sur le graphe et la dimension spectrale.

Le terme de courbure discrète peut être écrit sous la forme
\[
S_{\rm curv}[W]
=
\sum_{(i,j)}
W_{ij}\kappa_{ij},
\]
où \(\kappa_{ij}\) est la courbure d'Ollivier--Ricci de l'arête \((i,j)\).

Le terme de localité pénalise les connexions fortement non locales dans la géométrie émergente :
\[
S_{\rm loc}[W]
=
\lambda_{\rm loc}
\sum_{i,j}
W_{ij}
\left(d_{ij}^{\rm ent}[W]\right)^2.
\]
Ce terme exprime la tension centrale de BuP : assez d'intrication est nécessaire pour produire une géométrie connectée, mais trop d'intrication non locale détruit la notion de proximité spatiale.

Enfin, le terme topologique
\[
S_{\rm topo}[W]
\]
doit permettre l'existence d'excitations stables protégées par des invariants spectraux ou topologiques. Une forme possible est
\[
S_{\rm topo}[W]
=
\sum_{\gamma}
c_\gamma \Phi_\gamma[W],
\]
où \(\gamma\) parcourt certains cycles du graphe.

L'équation d'Euler--Lagrange du réseau est alors
\[
\frac{\delta S_{\rm spec}}{\delta W_{ij}}
+
\frac{\delta S_{\rm curv}}{\delta W_{ij}}
+
\frac{\delta S_{\rm loc}}{\delta W_{ij}}
+
\frac{\delta S_{\rm topo}}{\delta W_{ij}}
=
0.
\]

Cette équation n'est ni une équation pour une métrique fondamentale, ni une équation pour un champ de matière fondamental. C'est une équation de cohérence informationnelle du réseau.

% ─────────────────────────────────────────────────────────────
\section{Matière comme excitation stable du réseau}
\label{sec:matter-excitations}
% ─────────────────────────────────────────────────────────────

Dans une formulation pure BuP, la matière apparaît comme une classe particulière de perturbations localisées du réseau :
\[
W_{ij}
=
W_{ij}^{(0)}
+
\delta W_{ij}^{\rm loc}.
\]

Pour qu'une telle perturbation soit interprétable comme matière, elle doit satisfaire plusieurs conditions.

\paragraph{Localisation.}
La perturbation doit être concentrée dans une région finie de la géométrie émergente :
\[
\delta W_{ij}^{\rm loc}\simeq0
\]
lorsque les nœuds \(i,j\) sont éloignés de la région excitée.

\paragraph{Stabilité.}
Elle doit persister sous une dynamique effective du réseau, par exemple une descente de gradient sur \(S_{\rm BuP}[W]\).

\paragraph{Transport.}
Elle doit pouvoir se déplacer dans la géométrie émergente tout en conservant son identité.

\paragraph{Charge ou invariant.}
Elle doit porter une quantité conservée ou quasi conservée :
\[
Q[\delta W^{\rm loc}].
\]

\paragraph{Réponse géométrique.}
Elle doit modifier localement la dimension spectrale, la courbure ou la métrique émergente.

Dans la limite continue, une famille d'excitations localisées peut être décrite par des champs effectifs :
\[
\delta W_{ij}^{(a)}
\longrightarrow
\phi_a(x).
\]
Les champs de matière ne sont donc pas fondamentaux. Ils sont des variables collectives décrivant des modes stables du réseau.

Le tenseur énergie-impulsion effectif est alors écrit formellement :
\[
T_{\mu\nu}^{\rm matter}
=
\mathcal{T}_{\mu\nu}[\delta W^{\rm loc}].
\]
Une approximation schématique est
\[
T_{\mu\nu}^{\rm matter}
\sim
\left\langle
\delta W^{\rm loc}\delta W^{\rm loc}
\right\rangle_{\rm coarse}.
\]

% ─────────────────────────────────────────────────────────────
\section{Limites continues : géométrie et matière}
\label{sec:continuum-limits}
% ─────────────────────────────────────────────────────────────

La formulation pure BuP doit reproduire la physique effective dans plusieurs limites.

Dans la limite géométrique, les distances d'intrication
\[
d_{ij}^{\rm ent}
\]
définissent une métrique continue
\[
g_{\mu\nu}^{\rm ent}(x).
\]
Le Laplacien du graphe converge vers un opérateur de type Laplace--Beltrami :
\[
L
\longrightarrow
-\Delta_{g_{\rm ent}}.
\]
La dimension spectrale locale devient
\[
d_s(x)
=
-2\frac{d\log P_x(\tau)}{d\log\tau}.
\]

Dans la limite gravitationnelle, les variations lentes du réseau produisent une équation effective :
\[
G_{\mu\nu}[g^{\rm ent}]
=
8\pi G_{\rm eff}(d_s)T_{\mu\nu}^{\rm matter}
+
8\pi G T_{\mu\nu}^{\rm ent}[d_s].
\]

Dans la limite de matière, les perturbations localisées
\[
\delta W^{\rm loc}
\]
deviennent des champs effectifs :
\[
\delta W^{\rm loc}
\longrightarrow
\phi_a(x)
\longrightarrow
T_{\mu\nu}^{\rm matter}.
\]

Le programme unifié est donc :
\[
\frac{\delta S_{\rm BuP}[W]}{\delta W_{ij}}=0
\quad\Longrightarrow\quad
\begin{cases}
W_{ij}\rightarrow g_{\mu\nu}^{\rm ent}
\rightarrow G_{\mu\nu}[g^{\rm ent}],\\[4pt]
W_{ij}\rightarrow d_s(x)
\rightarrow G_{\rm eff}(d_s),\\[4pt]
\delta W^{\rm loc}\rightarrow T_{\mu\nu}^{\rm matter}.
\end{cases}
\]

% ─────────────────────────────────────────────────────────────
\section{Programme numérique}
\label{sec:numerical-program}
% ─────────────────────────────────────────────────────────────

Le test numérique central consiste à partir d'une matrice de fond
\[
W_{ij}^{(0)}
\]
et à injecter une perturbation localisée :
\[
W_{ij}
=
W_{ij}^{(0)}
+
A f_i f_j,
\]
où
\[
f_i
=
\exp\left[
-\frac{d_{\rm ent}^2(i,i_0)}{2\sigma^2}
\right].
\]

On définit une densité d'excitation nodale :
\[
\rho_i^{\rm exc}
=
\sum_j
|W_{ij}-W_{ij}^{(0)}|.
\]
On définit également un proxy de matière sur les arêtes :
\[
T_{ij}^{\rm matter,proxy}
=
(\delta W_{ij})^2.
\]

La réponse géométrique est mesurée par la variation du proxy d'Einstein discret :
\[
\delta G_{ij}^{\rm proxy}
=
G_{ij}^{\rm proxy}[W]
-
G_{ij}^{\rm proxy}[W^{(0)}],
\]
avec
\[
G_{ij}^{\rm proxy}
=
-\kappa_{ij}
+
\frac12R_{\rm edge}.
\]

L'observable principale est
\[
\mathcal{O}_{\rm matter}(\lambda)
=
\rho_{\rm Spearman}
\left(
T_{ij}^{\rm matter,proxy},
|\delta G_{ij}^{\rm proxy}|
\right).
\]

Si
\[
\mathcal{O}_{\rm matter}(\lambda)>0,
\]
l'excitation localisée agit comme une source géométrique effective. Si
\[
\mathcal{O}_{\rm matter}(\lambda)<0,
\]
la réponse devient anti-corrélée et l'interprétation en source matérielle locale échoue.

% ─────────────────────────────────────────────────────────────
\section{Résultats numériques : seuil critique de localité/matière}
\label{sec:results}
% ─────────────────────────────────────────────────────────────

Nous avons d'abord scanné l'observable \(\mathcal{O}_{\rm matter}(\lambda)\) pour des graphes \(N=20\). Un scan large avait initialement indiqué un changement de signe entre \(\lambda=0.57\) et \(\lambda=0.71\). Un scan fin autour de cette région affine fortement l'estimation.

Les deux modes d'excitation, positif et négatif, donnent le même bracket :
\[
0.58<\lambda_c<0.60.
\]
Par interpolation linéaire, on obtient :
\[
\lambda_c^{(+)}
\simeq0.5932,
\qquad
\lambda_c^{(-)}
\simeq0.5932.
\]
Nous retenons donc :
\[
\boxed{
\lambda_c(N=20)\simeq0.599
}
\]
comme estimation de taille finie, en tenant compte des scans large et fin.

La transition est abrupte. L'observable d'ordre passe d'une valeur positive d'environ
\[
\mathcal{O}_{\rm matter}\sim+0.66
\]
à \(\lambda=0.58\), à une valeur négative d'environ
\[
\mathcal{O}_{\rm matter}\sim-0.34
\]
à \(\lambda=0.60\). Le saut est donc de l'ordre de
\[
\Delta \mathcal{O}_{\rm matter}\sim1
\]
sur un intervalle
\[
\Delta\lambda=0.02.
\]

Le fait que les modes positif et négatif donnent le même seuil montre que la transition ne dépend pas du signe de la perturbation. Elle est donc interprétée comme une propriété du régime global du réseau.

\begin{figure}[htbp]
\centering
\includegraphics[width=0.95\textwidth]{../figures/fig_lambda_c_fine_positive.png}
\caption{
Scan fin du candidat paramètre d'ordre matière/localité pour \(N=20\), mode positif. L'observable \(\mathcal{O}_{\rm matter}(\lambda)\) change brutalement de signe entre \(\lambda=0.58\) et \(\lambda=0.60\), donnant une estimation interpolée \(\lambda_c\simeq0.5932\). La transition sépare un régime local où l'excitation agit comme source géométrique effective d'un régime non local où la réponse devient anti-corrélée.
}
\label{fig:lambda-c-positive}
\end{figure}

\begin{figure}[htbp]
\centering
\includegraphics[width=0.95\textwidth]{../figures/fig_lambda_c_fine_negative.png}
\caption{
Même scan critique pour le mode négatif. Le seuil obtenu est compatible avec celui du mode positif, ce qui indique que \(\lambda_c\) est une propriété du régime informationnel du graphe et non du signe de l'excitation injectée.
}
\label{fig:lambda-c-negative}
\end{figure}

% ─────────────────────────────────────────────────────────────
\section{Analyse de taille finie préliminaire}
\label{sec:fss}
% ─────────────────────────────────────────────────────────────

Afin de tester si le changement de signe est un artefact de taille finie, nous avons répété le scan de
\[
\mathcal{O}_{\rm matter}(\lambda)
\]
pour
\[
N=9,12,16,20.
\]

La Figure~\ref{fig:fss-lambda-c} montre que la transition n'est pas résolue aux petites tailles. Pour \(N=9\), l'observable reste presque plate autour de
\[
\mathcal{O}_{\rm matter}\simeq0.4.
\]
Pour \(N=12\), elle reste positive mais commence à décroître à grand \(\lambda\). Pour \(N=16\), la décroissance devient plus visible, sans produire de changement de signe. Enfin, pour \(N=20\), un crossing net apparaît autour de
\[
\lambda_c(N=20)\simeq0.599.
\]

Cette progression avec \(N\) suggère que la transition matière/localité est arrondie aux petites tailles et devient résolue lorsque la taille du graphe est suffisante. Nous interprétons donc \(\lambda_c(N=20)\) comme une estimation de taille finie d'un seuil critique candidat, et non encore comme une constante universelle.

\begin{figure}[htbp]
\centering
\includegraphics[width=0.92\textwidth]{../figures/fig_FSS_lambda_c.png}
\caption{
Analyse de taille finie préliminaire du candidat paramètre d'ordre
\(\mathcal{O}_{\rm matter}(\lambda)\) pour \(N=9,12,16,20\).
Pour \(N=9\), l'observable reste presque plate et positive. Pour \(N=12\) et \(N=16\), elle décroît progressivement avec \(\lambda\), sans crossing. Pour \(N=20\), un changement de signe net apparaît autour de \(\lambda_c(N=20)\simeq0.599\). Cette évolution suggère que la transition matière/localité devient résolue seulement à taille suffisamment grande.
}
\label{fig:fss-lambda-c}
\end{figure}

% ─────────────────────────────────────────────────────────────
\section{Interprétation : localité, matière et décohérence géométrique}
\label{sec:discussion}
% ─────────────────────────────────────────────────────────────

Les résultats numériques suggèrent que \(\lambda\) agit comme un candidat paramètre de contrôle de la phase informationnelle du réseau. Lorsque
\[
\lambda<\lambda_c,
\]
le réseau supporte des excitations localisées qui produisent une réponse de courbure positivement corrélée. Dans ce régime, les excitations peuvent être interprétées comme des sources matérielles effectives.

Lorsque
\[
\lambda>\lambda_c,
\]
la réponse devient anti-corrélée. L'excitation est encore présente dans \(W_{ij}\), mais elle ne se traduit plus par une réponse géométrique locale positive. Cela suggère une perte de cohérence géométrique locale.

Il est donc naturel d'interpréter \(\lambda_c\) comme un seuil de transition localité/décohérence. En dessous du seuil, la perturbation reste géométriquement cohérente. Au-dessus, elle devient dispersée ou réorganisée à l'échelle du réseau, et ne peut plus être interprétée comme source matérielle locale.

Cette lecture est plus naturelle que l'interprétation en antimatière. Les tests d'excitations positives et négatives montrent en effet que le changement de signe de \(\mathcal{O}_{\rm matter}\) ne dépend pas du signe de \(\delta W\). Les deux modes donnent le même seuil. Il s'agit donc d'une propriété du régime global du graphe, non d'un mode conjugué simple.

La conclusion conceptuelle est :
\[
\boxed{
\text{la matière émergente nécessite une phase locale du réseau d'intrication.}
}
\]

Cela donne une condition nouvelle à la formulation pure BuP. Il ne suffit pas qu'un réseau soit intriqué pour produire de la matière effective. Il doit être dans une phase où les excitations localisées restent géométriquement cohérentes.

% ─────────────────────────────────────────────────────────────
\section{Limites et perspectives}
% ─────────────────────────────────────────────────────────────

Cette étude reste préliminaire. Plusieurs limites doivent être soulignées.

Premièrement, le seuil
\[
\lambda_c(N=20)\simeq0.599
\]
est une estimation de taille finie. Des simulations à \(N>20\) seront nécessaires pour déterminer si
\[
\lambda_c(N)
\]
converge vers une valeur limite.

Deuxièmement, le proxy de matière utilisé ici,
\[
T_{ij}^{\rm matter,proxy}=(\delta W_{ij})^2,
\]
est volontairement minimal. Il mesure l'énergie informationnelle locale de la perturbation, mais ne capture pas encore des propriétés plus riches comme les charges, les phases ou les structures topologiques.

Troisièmement, la dynamique de stabilité temporelle n'est pas encore simulée. Le présent papier teste la réponse géométrique instantanée à une excitation localisée. Une étape ultérieure devra introduire une évolution effective du réseau :
\[
\frac{dW_{ij}}{dt_{\rm eff}}
=
-\eta
\frac{\delta S_{\rm BuP}}{\delta W_{ij}},
\]
puis vérifier la stabilité, la propagation et la conservation d'invariants.

Quatrièmement, les liens avec les équations de champ connues restent à construire. Ce papier ne dérive pas encore les équations de Dirac, Maxwell ou Yang--Mills. Il établit seulement une première condition numérique pour que des excitations du réseau puissent être interprétées comme sources matérielles effectives.

Les perspectives immédiates sont donc :
\begin{enumerate}
  \item tester \(N>20\) lorsque le coût numérique le permettra ;
  \item raffiner l'analyse de taille finie ;
  \item tester d'autres profils d'excitation ;
  \item introduire des perturbations topologiques ou orientées ;
  \item définir des charges informationnelles conservées ;
  \item simuler la dynamique temporelle des excitations.
\end{enumerate}

% ─────────────────────────────────────────────────────────────
\section{Conclusion}
\label{sec:conclusion}
% ─────────────────────────────────────────────────────────────

Ce papier propose une étape vers une formulation pure BuP, dans laquelle la matière et la géométrie émergent du même réseau d'intrication. L'équation fondamentale conjecturale est
\[
\frac{\delta S_{\rm BuP}[W]}{\delta W_{ij}}=0.
\]

Dans cette formulation, l'équation d'Einstein effective de Paper~5 n'est pas fondamentale. Elle apparaît comme une limite continue et coarse-grainée de la dynamique du réseau. De même, le tenseur énergie-impulsion de la matière n'est pas postulé : il doit émerger d'excitations localisées de \(W_{ij}\).

Le résultat numérique principal est la définition d'un candidat paramètre d'ordre
\[
\mathcal{O}_{\rm matter}(\lambda)
=
\rho_{\rm Spearman}
\left(
T_{ij}^{\rm matter,proxy},
|\delta G_{ij}^{\rm proxy}|
\right),
\]
qui mesure si une excitation localisée agit comme source géométrique effective.

Pour \(N=20\), cette observable change brutalement de signe autour de
\[
\lambda_c(N=20)\simeq0.599.
\]
Une analyse de taille finie préliminaire montre que la transition n'est pas visible à \(N=9\), reste arrondie à \(N=12\) et \(N=16\), puis devient résolue à \(N=20\).

Nous interprétons ce seuil comme un candidat seuil critique de localité/matière. En dessous de \(\lambda_c\), le réseau supporte des excitations localisées interprétables comme matière effective. Au-dessus, la réponse devient anti-corrélée, suggérant une perte de cohérence géométrique locale.

Ainsi, Paper~6 propose la thèse suivante :
\[
\boxed{
\text{la matière est une excitation stable d'une phase locale du réseau d'intrication.}
}
\]

Cette thèse reste à confirmer par des simulations plus grandes et par une dynamique explicite du réseau. Néanmoins, les résultats présentés ici fournissent un premier test numérique interne de l'idée que la matière, comme la géométrie, peut émerger du substrat informationnel BuP.

% ─────────────────────────────────────────────────────────────
\begin{thebibliography}{9}
% ─────────────────────────────────────────────────────────────

\bibitem{hamdad-paper2}
F.~Hamdad,
\textit{Cosmological constraints on an emergent spatial dimension in Bottom-Up Quantum Gravity},
BuP Paper 2, HAL hal-05590614v1, 2026.
\url{https://hal.science/hal-05590614v1}

\bibitem{hamdad-paper3}
F.~Hamdad,
\textit{Coherent effective gravitational coupling and the \(\sigma_8\) tension in Bottom-Up Quantum Gravity},
BuP Paper 3, GitHub repository, 2026.

\bibitem{hamdad-paper4}
F.~Hamdad,
\textit{Local dimensional reduction and early massive galaxies in Bottom-Up Quantum Gravity},
BuP Paper 4, GitHub repository, 2026.

\bibitem{hamdad-paper5}
F.~Hamdad,
\textit{Équation d'Einstein effective émergente dans la Gravité Quantique Bottom-Up},
BuP Paper 5, GitHub repository, 2026.
\url{https://github.com/Farid-Hamdad/Bottom-up-Quantum-Gravity/tree/bup_cosmology/paper5_effective_einstein}

\bibitem{ollivier}
Y.~Ollivier,
\textit{Ricci curvature of Markov chains on metric spaces},
Journal of Functional Analysis, 256, 810--864, 2009.

\bibitem{spectral-dimension}
D.~Benedetti,
\textit{Fractal properties of quantum spacetime},
Physical Review Letters, 102, 111303, 2009.

\end{thebibliography}

\end{document}
